\documentclass[12pt]{article}
\usepackage[utf8]{inputenc}
\usepackage[russian]{babel}
\usepackage{amsmath}

\begin{document}

\section*{Замечательные пределы}

\begin{enumerate}
    \item Первый замечательный предел:
    \[
    \lim_{x \to 0} \frac{\sin x}{x} = 1
    \]
    
    \item Второй замечательный предел (первая форма):
    \[
    \lim_{x \to 0} (1 + x)^{\frac{1}{x}} = e
    \]
    
    \item Второй замечательный предел (вторая форма):
    \[
    \lim_{x \to \infty} \left(1 + \frac{1}{x}\right)^x = e
    \]
\end{enumerate}

\section*{Правило Лопиталя}

Если при вычислении предела предела функций получается неопределенность вида $\left[\frac{0}{0}\right]$ или $\left[\frac{\infty}{\infty}\right]$, то предел отношения функций равен пределу отношения их производных:
\[
\lim_{x \to a} \frac{f(x)}{g(x)} = \lim_{x \to a} \frac{f'(x)}{g'(x)}
\]

\end{document}